\begingroup
\renewcommand{\arraystretch}{1.45}
\newcommand{\TransscendentArticle}[1]{\par\medskip\Needspace{5\baselineskip}\pdfbookmark[1]{Art. #1}{transscendentes-ii-art-#1}\centerline{[#1.]}\nobreak\smallskip}
\clearpage\selectlanguage{ngerman}\fancyhead[L]{\small ZUR THEORIE DER NEUEN TRANSZENDENTEN. II.}\pdfbookmark[0]{Zur Theorie der transzendenten Funktionen}{transscendentes-ii}\SourceLeaf{445}
\begin{center}[II.]\end{center}
\begin{center}\textbf{ZUR THEORIE DER TRANSZENDENTEN FUNKTIONEN GEHÖRIG.}\end{center}
\par\nobreak\medskip\centerline{\rule{.14\linewidth}{.4pt}}
\TransscendentArticle{1}
Es sei $\sum e^{-\alpha(k+\omega)^2}=T$, wobei für $k$ alle ganzen positiven und negativen
Zahlen eingesetzt werden, und
\[
T=A+2B\cos\omega P+2C\cos2\omega P+2D\cos3\omega P+\text{ usw.}
\]
\Needspace{15\baselineskip}
so gilt
\[
\begin{gathered}
A=\int T\,\mathrm{d}\omega,\qquad B=\int T\cos\omega P\mathbin{.}\mathrm{d}\omega,\qquad
\\[7pt]\quad{}C=\int T\cos2\omega P\mathbin{.}\mathrm{d}\omega,\\[7pt]
D=\int T\cos3\omega P\mathbin{.}\mathrm{d}\omega,\quad\text{usw.}
\end{gathered}
\]
\Needspace{18\baselineskip}
wobei alle Integrale von $\omega=0$ bis $\omega=1$ ausgedehnt werden. Es ist jedoch klar, dass jene Integrale zwischen diesen Grenzen mit den Integralen
\[
\begin{gathered}\int e^{-\alpha\omega\omega}\,\mathrm{d}\omega,\quad
\int e^{-\alpha\omega\omega}\cos\omega P\mathbin{.}\mathrm{d}\omega,\quad
\\[7pt]\quad{}\int e^{-\alpha\omega\omega}\cos2\omega P\mathbin{.}\mathrm{d}\omega,\quad
\int e^{-\alpha\omega\omega}\cos3\omega P\mathbin{.}\mathrm{d}\omega,\text{ usw.}\end{gathered}
\]
\Needspace{20\baselineskip}
übereinstimmen, wenn diese von $\omega=-\infty$ bis $\omega=+\infty$ ausgedehnt werden. Da nun allgemein
\[
\begin{gathered}
e^{-\alpha\omega\omega}\cos n\omega P\\[7pt]
=\tfrac12e^{-\alpha\omega\omega+in\omega P}+\tfrac12e^{-\alpha\omega\omega-in\omega P}
\\[7pt]\quad{}=\tfrac12e^{-\frac{nnPP}{4\alpha}}
\Bigl\{e^{-\alpha\Bigl(\omega-\frac{inP}{2\alpha}\Bigr)^2}
+e^{-\alpha\Bigl(\omega+\frac{inP}{2\alpha}\Bigr)^2}\Bigr\}
\end{gathered}
\]


\Needspace{8\baselineskip}
\SourceLeaf{446}
so folgt leicht, dass
\[
\int e^{-\alpha\omega\omega}\cos n\omega P\mathbin{.}\mathrm{d}\omega
=e^{-\frac{nn\pi\pi}{\alpha}}\surd\frac{\pi}{\alpha}
\]
\Needspace{12\baselineskip}
folglich
\[
\begin{gathered}T=\surd\frac{\pi}{\alpha}\mathbin{.}
\Bigl\{1+2e^{-\frac{\pi\pi}{\alpha}}\cos\omega P
\\[7pt]\quad{}+2e^{-4\frac{\pi\pi}{\alpha}}\cos2\omega P
+2e^{-9\frac{\pi\pi}{\alpha}}\cos3\omega P+\text{ usw.}\Bigr\}\end{gathered}
\]
\Needspace{9\baselineskip}
In einer anderen Form:
\[
\sum e^{-\alpha(k+\omega)^2}
=\surd\frac{\pi}{\alpha}\mathbin{.}e^{-\alpha\omega\omega}\mathbin{.}
\sum e^{-\frac{\pi\pi}{\alpha}\left(k+\frac{\alpha\omega i}{\pi}\right)^2}
\]
\Needspace{15\baselineskip}
Allgemeiner:
\[
\begin{gathered}
\sum e^{-\alpha(k+\omega)^2}\\[7pt]\quad{}\{\cos2(k+\omega)\alpha\psi-i\sin2(k+\omega)\alpha\psi\}\\[7pt]
=\surd\frac{\pi}{\alpha}\mathbin{.}
\sum e^{-\frac{\pi\pi}{\alpha}\Bigl(k-\frac{\alpha\psi}{\pi}\Bigr)^2}
\{\cos2k\omega\pi-i\sin2k\omega\pi\}
\end{gathered}
\]
\Needspace{17\baselineskip}
Oder, setzt man $\alpha\alpha'=\pi\pi$ und $-\alpha\psi=\omega'\pi$,
\[
\begin{gathered}
\sum e^{-\alpha(k+\omega)^2}\\[7pt]\quad{}\{\cos(2k+2\omega)\omega'\pi-i\sin(2k+2\omega)\omega'\pi\}\\[7pt]
=\surd\frac{\pi}{\alpha}\mathbin{.}\sum e^{-\alpha'(k+\omega')^2}
\{\cos2k\omega\pi-i\sin2k\omega\pi\}
\end{gathered}
\]
\[
\begin{gathered}\sum e^{-\alpha(k+\omega-\frac14)^2}-\sum e^{-\alpha(k+\omega+\frac14)^2}
\\[7pt]\quad{}=4\surd\frac{\pi}{\alpha}\mathbin{.}
\Bigl\{e^{-\frac{\pi\pi}{\alpha}}\sin\omega P
-3e^{-9\frac{\pi\pi}{\alpha}}\sin3\omega P+\text{ usw.}\Bigr\}\end{gathered}
\]
\Needspace{9\baselineskip}
Man kann den Lehrsatz auch so ausdrücken:
\[
\sum t e^{-\pi t^{4}kk-2\pi ttk\surd u-\frac12\pi u}
\]
ändert seinen Wert nicht, wenn $t$ in $\dfrac1t$ und $u$ in $-u$ umgewandelt wird.
\clearpage
\endgroup
